\begin{trapbox}

	\begin{description}[style=nextline, leftmargin=2.8em, labelsep=0.5em, itemsep=0.35em]
		\item[\textbf{\textcolor{CTrap}{易错点一（位置混淆）}}]
			误认为极线总是切线，当点 $P$ 在曲线外时，也将极线当作切线使用。
		\item[\textbf{\textcolor{CTrap}{正确理解（区分内外）：}}]
			点 $P$ 在曲线上时极线是切线；点 $P$ 在曲线外时极线是切点弦（连接两切点的直线）。
		\item[\textbf{\textcolor{CTrap}{错因分析（忽视位置）：}}]
			对极线的定义理解不完整，未根据点与曲线的位置关系区分几何意义。
	\end{description}

	\medskip
	\begin{description}[style=nextline, leftmargin=2.8em, labelsep=0.5em, itemsep=0.35em]
		\item[\textbf{\textcolor{CTrap}{易错点二（符号混淆）}}]
			在写双曲线的极线方程时，误写为加号，即 $\dfrac{x_0x}{a^2}+\dfrac{y_0y}{b^2}=1$。
		\item[\textbf{\textcolor{CTrap}{正确理解（符号匹配）：}}]
			双曲线极线方程应为 $\dfrac{x_0x}{a^2}-\dfrac{y_0y}{b^2}=1$，中间为减号。
		\item[\textbf{\textcolor{CTrap}{错因分析（机械记忆）：}}]
			只记忆了椭圆极线方程的形式，未注意双曲线标准方程的符号差异。
	\end{description}

	\medskip
	\begin{description}[style=nextline, leftmargin=2.8em, labelsep=0.5em, itemsep=0.35em]
		\item[\textbf{\textcolor{CTrap}{易错点三（系数错误）}}]
			对抛物线 $y^2=2px$，误将极线方程写成 $y_0y=2p(x+x_0)$ 或 $y_0y=p(x_0x)$。
		\item[\textbf{\textcolor{CTrap}{正确理解（系数对应）：}}]
			抛物线 $y^2=2px$ 的极线方程为 $y_0y=p(x+x_0)$，右侧一次项系数直接为 $p$，括号内为 $x+x_0$。
		\item[\textbf{\textcolor{CTrap}{错因分析（记忆偏差）：}}]
			混淆了抛物线标准形式中 $2p$ 与极线方程中 $p$ 的关系，或未正确替换 $x$。
	\end{description}

\end{trapbox}
